\subsection{The precise incidence quotient and the remaining regular-quotient problem}
\paragraph{J1i (Known construction, with the family-kernel consequence proved here).}
Let $k$ be an algebraically closed field of any characteristic, let
$X/k$ be smooth projective of pure dimension $n$, and write
$A^i(X)=\operatorname{CH}^i(X)_{\rm alg}$. Define
\[
 I^i(X)=\bigcap_{T,w}\ker\left(
 w_*:A^i(X)\longrightarrow\operatorname{Pic}^0(T)(k)\right),
 \qquad w\in\operatorname{CH}^{n-i+1}(X\times T),
\]
where $T$ runs through smooth projective $k$-varieties; on $k$-points
$\operatorname{Pic}^0(T)$ means its reduced identity component.
This is \emph{incidence equivalence}, also called Picard equivalence
in the older sources. It is not a definition by all regular
homomorphisms to abelian varieties.

Kleiman explicitly separates the two relations in
\cite[pp.~446--447]{DGJ_Kleiman1971}. His p.~446 attributes the
boundedness argument for incidence equivalence to an unpublished
1964 observation of Grothendieck, and p.~447, footnote~5, records
the complementary duality with finite kernel. Thus the historical
connection with 44:2--9 has a primary-source basis, without
identifying two different equivalence relations.
The complete construction is Saito's
\cite[Definition 3.1, p.~108; Proposition 3.4, pp.~109--110;
Theorem 3.6, pp.~110--111]{DGJ_Saito1979}: there is an abelian
variety $P^i_{\rm inc}(X)$ and a surjective regular Picard
homomorphism
\[
 \theta_X^i:A^i(X)\longrightarrow P^i_{\rm inc}(X)(k),
 \qquad\ker\theta_X^i=I^i(X),
 \qquad 2\dim P^i_{\rm inc}(X)\leq b^{2i-1}(X).
\]
Here the Betti number is the dimension of $\ell$-adic cohomology
for any prime $\ell\ne\operatorname{char}k$. The map is universal
among Picard homomorphisms, which are the homomorphisms obtained
by a complementary correspondence to the Picard variety of a
smooth projective variety, with the image viewed as an abelian
subvariety. All these statements hold in every codimension,
including characteristics two and three. Saito's bound is written
using $b^{2n-2i+1}$; Poincaré duality gives the displayed index.

Here is the exact family-kernel statement which follows. Let
$(S,s_0)$ be smooth projective and pointed, and let
$\xi\in\operatorname{CH}^i(S\times X)$. Correspondence
functoriality supplies a homomorphism of $k$-abelian varieties
\[
 f_\xi:\operatorname{Alb}(S)\longrightarrow P^i_{\rm inc}(X),
 \qquad
 f_\xi([s]-[s_0])=\theta_X^i(\xi_s-\xi_{s_0}).
\]
For each $w$ in the definition of $I^i(X)$, composition with
$\xi$ is a divisor correspondence on $S\times T$ and defines
$g_{w,\xi}:\operatorname{Alb}(S)\to\operatorname{Pic}^{\rm var}(T)$.
Define $N_{\rm inc}(\xi)$ as the scheme-theoretic intersection of
all their kernels. Then
\[
 N_{\rm inc}(\xi)=\ker f_\xi,
 \qquad
 \operatorname{Alb}(S)/N_{\rm inc}(\xi)
     \simeq\operatorname{Im}(f_\xi),
 \qquad
 \dim\operatorname{Alb}(S)/N_{\rm inc}(\xi)
       \leq\tfrac12 b^{2i-1}(X).                 \tag{J1i.1}
\]
The quotient and the image here are abelian varieties; the kernel
may be nonreduced. To prove the equality, universality makes every
$g_{w,\xi}$ factor through $f_\xi$, giving one inclusion.
Conversely, $\theta_X^i$ is itself a Picard homomorphism
\cite[(3.5), p.~110]{DGJ_Saito1979}. Thus for some $T_0,w_0$,
its composition with a closed immersion
$P^i_{\rm inc}(X)\hookrightarrow\operatorname{Pic}^{\rm var}(T_0)$
is $(w_0)_*$. The corresponding $g_{w_0,\xi}$ has kernel exactly
$\ker f_\xi$, giving the other inclusion as group schemes, not
only on geometric points. The fppf quotient by a homomorphism's
kernel is its abelian-variety image. The dimension assertion now
follows from Saito's bound. This proves (J1i.1).

Saito further constructs, using his explicitly specified
miniversal-cycle integers, an isogeny
\[
 P^i_{\rm inc}(X)\longrightarrow
          P^{n+1-i}_{\rm inc}(X)^\vee;
\]
see \cite[(4.1)--(4.2), p.~112; Proposition 4.4 and
Corollary 4.5, p.~113]{DGJ_Saito1979}. This is a known
\emph{isogeny} statement. Its integer factors must be retained;
it does not supply the unscaled integral intersection isomorphism
asserted in J.2. Therefore the incidence reading of the
representability and boundedness portions of Theorems~1--3 is
settled by the cited construction and (J1i.1). The stronger
integral duality and the identification with a quotient defined
by all regular maps are separate assertions.

\paragraph{The regular interpretation and its exact remaining gate.}
A homomorphism $h:A^i(X)\to B(k)$ to a $k$-abelian variety is
\emph{regular} when every pointed smooth projective family of
cycles induces a morphism to $B$. Put
\[
 I^i_{\rm ab}(X)=\bigcap_h\ker h.
\]
This is Samuel's abelian equivalence. Every Picard homomorphism is
regular, so
$I^i_{\rm ab}(X)\subseteq I^i(X)$. A universal regular
homomorphism, if it exists, therefore has a canonical surjective
homomorphism from its abelian variety to $P^i_{\rm inc}(X)$.
This observation does not prove that the kernel is finite or zero.

For a fixed family $\xi$, let $N_{\rm ab}(\xi)$ be the
scheme-theoretic intersection of the kernels of all homomorphisms
$h_\xi:\operatorname{Alb}(S)\to B$ obtained from regular $h$.
Then the induced group homomorphism
\[
 \operatorname{Alb}(S)(k)\longrightarrow
                    A^i(X)/I^i_{\rm ab}(X)
\]
has kernel $N_{\rm ab}(\xi)(k)$, without assuming a universal
regular representative. Indeed, every composite on degree-zero
zero-cycles of $S$ is regular and factors through its Albanese;
hence the Albanese kernel is killed modulo $I^i_{\rm ab}(X)$.
The displayed map is well defined because the Albanese map on
zero-cycles is surjective over an algebraically closed field, and
its kernel is the stated intersection by definition. The
intersection is a closed subgroup scheme: on a noetherian
scheme the sum of the defining ideal sheaves is generated by a
finite subfamily. Thus each individual family has an abelian
quotient. What is missing is a uniform bound as the family varies.

More precisely, the existence of such a uniform dimension bound
is equivalent to existence of a universal regular homomorphism.
The bounded-target criterion is
\cite[Theorem 2.2(ii), p.~103 and proof pp.~103--106]{DGJ_Saito1979}
for the class of all regular homomorphisms, identified in his
Example (1.3.1), p.~101. To apply it to the family formulation,
any surjective regular $h$ admits a family on an abelian variety
whose induced map to its target is surjective
\cite[Proposition 1.2, pp.~98--100]{DGJ_Saito1979}; this target
is a quotient of the corresponding family quotient, so a uniform
family bound bounds all regular targets. Conversely, a universal
regular representative is itself among the tests defining every
$N_{\rm ab}(\xi)$, so each such quotient is its abelian
subvariety image. This proves the claimed equivalence of the two
boundedness problems. A bound by the particular Betti number is
stronger than mere existence.

The known existence cases are $i=1,2,n$.
For codimension two the complete repair of Murre's proof is
\cite[Theorem 1 and its two proofs, author version
pp.~2--4]{DGJ_Kahn2021}; for all three indices the Betti bound is
\cite[Proposition 8.2, p.~31]{DGJ_ACMVFunctor}.
That proposition explicitly restricts its cohomological injection
to $i=1,2,n$. Applying it in arbitrary codimension would assume
the missing injectivity of the higher Bloch map and would falsely
close the remaining gate.

A concrete unresolved test is Kahn's Conjecture~1, author version
p.~5: take a generic abelian threefold $A$ defined over a field
$k_0$ finitely generated over $\mathbb Q$, a nonisotrivial elliptic
curve $E/k_0(t)$, and
\[
 V=(A_{k_0(t)}\times E)_{\overline{k_0(t)}}.
\]
The conjecture states that a universal regular homomorphism on
$A^3(V)$ does not exist. Its preceding infinite-torsion result
alone does not prove this conjecture. Although every surjective
regular homomorphism is surjective on prime-to-characteristic
primary torsion \cite[Remark 2, p.~4]{DGJ_Kahn2021}, an infinite
torsion kernel may be killed in every finite-dimensional target.
To refute universal existence one must still construct regular
quotients of unbounded dimension, or another actual obstruction
satisfying the regularity condition. Neither the torsion argument
nor the opened motivic constructions provides that step here.
Accordingly the all-regular version in codimension three on
fourfolds, and its requested general form, remain unresolved.

\paragraph{A single self-incidence kernel is not the common kernel.}
The letter at 44:9 constructs $N$ from one symmetric
self-correspondence on a parameter variety. Such a kernel cannot
be substituted for $N_{\rm inc}(\xi)$ without an additional
argument. The distinction is visible in the same explicit
threefold as J2b. Let $X=E\times\mathbb P^1\times\mathbb P^1$,
let $S=E^\vee$, and set $\xi=c_1(\mathcal P)h_1$, where
$\mathcal P$ is the normalized Poincaré bundle. Its
self-incidence divisor is zero because $h_1^2=0$, so its
self-kernel is all of $S$. But pairing with the independent
family $c_1(\mathcal P)h_2$ gives the elliptic duality
$S\to S^\vee$, up to sign, which is an isomorphism. Hence
$N_{\rm inc}(\xi)=0$. Equivalently the family map is the
inclusion of the first elliptic factor in $P^2_{\rm inc}(X)=S^2$.
The projective-bundle decomposition and the two complementary
Poincaré tests prove this identification integrally, just as in
J2b. Thus one self-pairing can miss an entire abelian family that
other incidence tests detect. This refutes that particular
shortcut from 44:9 to Theorem~1; it does not decide an undefined
$N(\xi)$ on 44:2 by assigning it a new definition.

\paragraph{Exact status of the remaining Lefschetz assertion.}
Theorem~6 is still unresolved in its full stated generality.
Even after choosing the incidence quotient, its three conclusions
are assertions about algebraic correspondences and abelian
varieties. Cohomological weak Lefschetz identifies the ambient
cohomology in its stable range, but does not show that an inverse
on the required algebraic part is induced by correspondences.
For the universal regular quotient, general existence is already
the preceding gate. For the literal manuscript, the two
transcriptions additionally disagree between an isomorphism and
an isogeny in the stable range. Neither a cohomological
isomorphism nor Saito's complementary isogeny proves the stronger
integral restriction claim. These are retained as explicit
unresolved assertions, rather than being replaced with a
cohomological theorem.
